% ===== I-1 (multiple_choice, difficulty=1) =====
\begin{question}[type=multiple_choice,difficulty=1,code={I-1}]
Cho hàm số $y = f(x)$ có đồ thị như hình vẽ. Tiệm cận ngang của đồ thị hàm số là đường thẳng có phương trình:

\image[alt={Hình minh hoạ}]{https://pending-r2.local/raw/01_a7c080c887.jpg}


\begin{choice}[label=A]
x = -1
\end{choice}

\begin{choice}[label=B]
x = 1
\end{choice}

\begin{choice}[correct,label=C]
y = 2
\end{choice}

\begin{choice}[label=D]
y = -2
\end{choice}

\end{question}

% ===== I-2 (multiple_choice, difficulty=1) =====
\begin{question}[type=multiple_choice,difficulty=1,code={I-2}]
Cho hàm số y = f(x) có bảng biến thiên như sau
Hàm số đã cho đồng biến trên khoảng nào dưới đây?

\image[alt={Hình minh hoạ}]{https://pending-r2.local/raw/03_a846cfb472.jpg}


\begin{choice}[label=A]
$ (-1;1) $
\end{choice}

\begin{choice}[label=B]
$ (-1;+\infty) $
\end{choice}

\begin{choice}[correct,label=C]
$ (0;1) $
\end{choice}

\begin{choice}[label=D]
$ (-1;0) $
\end{choice}

\end{question}

% ===== I-3 (multiple_choice, difficulty=1) =====
\begin{question}[type=multiple_choice,difficulty=1,code={I-3}]
Cho cấp số cộng $ (u_n) $ có số hàng đầu $ u_1 = 3 $ và công sai $ d = -4 $. Số hàng thứ năm của cấp số công là

\begin{choice}[label=A]
-3072
\end{choice}

\begin{choice}[correct,label=B]
-13
\end{choice}

\begin{choice}[label=C]
-17
\end{choice}

\begin{choice}[label=D]
768
\end{choice}

\end{question}

% ===== I-4 (multiple_choice, difficulty=2) =====
\begin{question}[type=multiple_choice,difficulty=2,code={I-4}]
Cho một chất điểm chuyển động có phương trình $s = 2t^2 + 3t$ (t tính bằng giây, s tính bằng mét). Vận tốc của chất điểm tại thời điểm $t_0 = 2$ (giây) bằng

\begin{choice}[label=A]
$19(m/s)$
\end{choice}

\begin{choice}[label=B]
$22(m/s)$
\end{choice}

\begin{choice}[correct,label=C]
$11(m/s)$
\end{choice}

\begin{choice}[label=D]
$9(m/s)$
\end{choice}

\end{question}

% ===== I-5 (multiple_choice, difficulty=2) =====
\begin{question}[type=multiple_choice,difficulty=2,code={I-5}]
Doanh thu bán hàng trong 20 ngày được lựa chọn ngẫu nhiên của một cửa hàng được ghi lại ở bảng sau (đơn vị: triệu đồng):

$$\begin{array}{|c|c|c|c|c|c|}
\hline
\text{Doanh thu} & \text{[5;7)} & \text{[7;9)} & \text{[9;11)} & \text{[11;13)} & \text{[13;15)} \\
\hline
\text{Số ngày} & \text{2} & \text{7} & \text{7} & \text{3} & \text{1} \\
\hline
\end{array}$$

Số trung bình của mẫu số liệu trên thuộc khoảng nào trong các khoảng dưới đây?

\begin{choice}[label=A]
[11;13)
\end{choice}

\begin{choice}[label=B]
[13;15)
\end{choice}

\begin{choice}[correct,label=C]
[9;11)
\end{choice}

\begin{choice}[label=D]
[7;9)
\end{choice}

\end{question}

% ===== I-6 (multiple_choice, difficulty=3) =====
\begin{question}[type=multiple_choice,difficulty=3,code={I-6}]
Cho hình chóp S.ABCD có đày ABCD là hình chữ nhất với AD = a, cạnh bên SẠ vuông góc với đáy và SA = a√3. Số đo của góc nhị diện [S, DC, B] bằng:

\begin{choice}[label=A]
$ 90^{\circ} $
\end{choice}

\begin{choice}[correct,label=B]
$ 60^{\circ} $
\end{choice}

\begin{choice}[label=C]
$ 30^{\circ} $
\end{choice}

\begin{choice}[label=D]
$ 45^{\circ} $
\end{choice}

\end{question}

% ===== I-7 (multiple_choice, difficulty=1) =====
\begin{question}[type=multiple_choice,difficulty=1,code={I-7}]
Cho hàm số y = f(x) có đồ thị như hình vẽ:
Khi đó, $ \lim_{x \to +\infty} f(x) $ bằng:

\image[alt={Hình minh hoạ}]{https://pending-r2.local/raw/07_e6942f4225.jpg}


\begin{choice}[label=A]
-\infty
\end{choice}

\begin{choice}[label=B]
+\infty
\end{choice}

\begin{choice}[label=C]
1
\end{choice}

\begin{choice}[correct,label=D]
-1
\end{choice}

\end{question}

% ===== I-8 (multiple_choice, difficulty=2) =====
\begin{question}[type=multiple_choice,difficulty=2,code={I-8}]
Nghiệm của phương trình $ \sqrt{3} + 3 \tan x = 0 $ là:

\begin{choice}[label=A]
$ x = \frac{\pi}{2} + k\pi $ ( $ k \in \mathbb{Z} $)
\end{choice}

\begin{choice}[correct,label=B]
$ x = -\frac{\pi}{6} + k\pi $ ( $ k \in \mathbb{Z} $)
\end{choice}

\begin{choice}[label=C]
$ x = \frac{\pi}{3} + k\pi $ ( $ k \in \mathbb{Z} $)
\end{choice}

\begin{choice}[label=D]
$ x = \frac{\pi}{2} + k2\pi $ ( $ k \in \mathbb{Z} $)
\end{choice}

\end{question}

% ===== I-9 (multiple_choice, difficulty=2) =====
\begin{question}[type=multiple_choice,difficulty=2,code={I-9}]
Cho hình lập phương $ ABCD \cdot A'B'C'D' $. Mệnh đề nào sau đây sai?

\begin{choice}[label=A]
$ \overrightarrow{AC} = \overrightarrow{AB} + \overrightarrow{AD} $
\end{choice}

\begin{choice}[correct,label=B]
$ \overrightarrow{AB} = \overrightarrow{CD} $
\end{choice}

\begin{choice}[label=C]
$ \left|\overrightarrow{AB}\right| = \left|\overrightarrow{CD}\right| $
\end{choice}

\begin{choice}[label=D]
$ \overrightarrow{AB} + \overrightarrow{AD} + \overrightarrow{AA'} = \overrightarrow{AC'} $
\end{choice}

\end{question}

% ===== I-10 (multiple_choice, difficulty=2) =====
\begin{question}[type=multiple_choice,difficulty=2,code={I-10}]
Trong không gian với hệ tọa độ Oxyz, cho hình bình hành ABCD có $ A(2;1,-3) $, $ B(0,-2;2,C(4,-3;0)) $. Toạ độ điểm D là

\begin{choice}[label=A]
$ (2;0,-5) $
\end{choice}

\begin{choice}[correct,label=B]
$ (6;0,-5) $
\end{choice}

\begin{choice}[label=C]
$ (2;0;5) $
\end{choice}

\begin{choice}[label=D]
$ (6;0;5) $
\end{choice}

\end{question}

% ===== I-11 (multiple_choice, difficulty=2) =====
\begin{question}[type=multiple_choice,difficulty=2,code={I-11}]
Tìm a đế hàm số $ \log_a x $ (0 < a $ \neq 1 $) có đồ thị là hình bên

\image[alt={Hình minh hoạ}]{https://pending-r2.local/raw/10_2d2b8074c5.jpg}


\begin{choice}[label=A]
$ a=2 $
\end{choice}

\begin{choice}[correct,label=B]
$ a=\sqrt{2} $
\end{choice}

\begin{choice}[label=C]
$ a=-\frac{1}{2} $
\end{choice}

\begin{choice}[label=D]
$ a=\frac{1}{2} $
\end{choice}

\end{question}

% ===== I-12 (multiple_choice, difficulty=1) =====
\begin{question}[type=multiple_choice,difficulty=1,code={I-12}]
Nghiệm của phương trình $ \log_{2} x = 3 $ là

\begin{choice}[label=A]
x = 2
\end{choice}

\begin{choice}[label=B]
x = 9
\end{choice}

\begin{choice}[correct,label=C]
x = 8
\end{choice}

\begin{choice}[label=D]
x = 3
\end{choice}

\end{question}

% ===== II-1 (true_false, difficulty=3) =====
\begin{question}[type=true_false,difficulty=3,code={II-1}]
Cho hàm số y = f(x) = ax³ + bx² + cx + d có đồ thị như hình vẽ dưới đây:

\image[alt={Hình minh hoạ}]{https://pending-r2.local/raw/12_dd89678b7d.jpg}


\begin{statement}[correct=false,label=a]
$ 2a + 3b + c = 9 $.
\end{statement}

\begin{statement}[correct=false,label=b]
Hàm số đạt cực tiêu tại x = 1.
\end{statement}

\begin{statement}[correct=true,label=c]
Tổng giá trị lớn nhất, giá trị nhỏ nhất của hàm số trên đoạn $ [-1;0] $ bằng 3.
\end{statement}

\begin{statement}[correct=false,label=d]
Hàm số đồng biến trên khoảng (-∞; -1).
\end{statement}

\end{question}

% ===== II-2 (true_false, difficulty=3) =====
\begin{question}[type=true_false,difficulty=3,code={II-2}]
Cho hình lăng trụ đứng ABC. A'B'C' có đáy ABC là tam giác vuông tại A. Gọi E, F lần lượt là trung điểm của AB và AA'. Cho biết AB = 2, BC = $ \sqrt{13} $, CC' = 4.

\begin{statement}[correct=false,label=a]
Thế tích khối lăng trụ ABC. A'B'C' bằng 8.
\end{statement}

\begin{statement}[correct=false,label=b]
Khoảng cách giữa hai đường thẳng A'C và FE bằng $ \frac{6}{7} $
\end{statement}

\begin{statement}[correct=true,label=c]
Đường thẳng AB vương góc với đường thẳng AC'
\end{statement}

\begin{statement}[correct=true,label=d]
Côn sân của góc giữa đường thẳng A'C và mặt phẳng đáy (ABC) bằng $ \frac{3}{5} $.
\end{statement}

\end{question}

% ===== II-3 (true_false, difficulty=3) =====
\begin{question}[type=true_false,difficulty=3,code={II-3}]
Hình về sau mô tả vị trí của máy bay vào thời điểm 9h30 phút. Biết các đơn vị trên hình tính theo đơn vị kim.

\image[alt={Hình minh hoạ}]{https://pending-r2.local/raw/14_6313b75802.jpg}


\begin{statement}[correct=true,label=a]
Phi công để máy bay ở chế độ tự động với vận tốc theo hướng động là 750km/h, độ cao không đổi. Biết rằng gió thổi theo hướng động với vận tốc 10m/s. Giả sử vận tốc và hướng gió không đổi thì lúc 10h30phút máy bay ở tọa độ (150;1086;9).
\end{statement}

\begin{statement}[correct=false,label=b]
Tọa độ của máy bay vào lúc 9h30 phút là: (300;150;9).
\end{statement}

\begin{statement}[correct=true,label=c]
Vào thời điểm 9h30 phút máy bay ở độ cao 9km.
\end{statement}

\begin{statement}[correct=false,label=d]
Sau khi bay đến vị trí lúc 10h30 thì máy bay ngược lại với vận tốc 800km/h với độ cao không đổi, biết lúc đó trời lặng gió thì lúc 11h máy bay ở toa độ (686;150;9).
\end{statement}

\end{question}

% ===== II-4 (true_false, difficulty=4) =====
\begin{question}[type=true_false,difficulty=4,code={II-4}]
Một doanh nghiệp kinh doanh sản xuất đồng hồ có đồ thị hàm tổng chi phí theo số sản phẩm, là một phần của đồ thị hàm số $ f(x)=\frac{ax^2+bx+c}{x+e} $ như hình vẽ (mỗi đơn vị trên trục hoành tương ứng với 100 sản phẩm, mỗi đơn vị trên trục tung tương ứng với 1000 USD). Biết rằng tầm đối xứng của đồ thị hàm số đó là điểm $ I\left(-1;\frac{2}{3}\right) $ và đường tiệm cận xiên của đồ thị đó đi qua điểm $ B(3;2) $.

\image[alt={Hình minh hoạ}]{https://pending-r2.local/raw/16_ee6adab79a.jpg}


\begin{statement}[correct=false,label=a]
Hàm số đồng biến trên khoảng $(0;+\infty)$.
\end{statement}

\begin{statement}[correct=false,label=b]
Đồ thị hàm số có đường tiệm cận đứng là $x = 1$.
\end{statement}

\begin{statement}[correct=true,label=c]
Hàm số có thể viết lại dưới dạng $f(x)=\frac{1}{3}x+1+\frac{d}{x+1}$, với $d \in \mathbb{R}$.
\end{statement}

\begin{statement}[correct=false,label=d]
Theo khảo sát, tổng doanh thu được mô tả bằng hàm số $R(x) = x^{2} + 2x$ và lợi nhuận khi bán 200 sản phẩm là 5250 USD. Khi chi phí nhỏ nhất, số sản phẩm (làm tròn) là 25 sản phẩm.
\end{statement}

\end{question}

% ===== III-1 (short_answer, difficulty=3) =====
\begin{question}[type=short_answer,difficulty=3,code={III-1}]
Cho hình chóp S.ABCD có đấy ABCD là hình thoi cạnh a, ADC = 60°, SA ⊥ (ABCD), SA = $ \sqrt{3}a $
. G là trọng tam tam giác SAC. Khoảng cách từ G đến (SCD) bằng: $ \frac{a\sqrt{m}}{n} $. Tính m+n?

\answer{30}

\end{question}

% ===== III-2 (short_answer, difficulty=1) =====
\begin{question}[type=short_answer,difficulty=1,code={III-2}]
Trọng lực $\overrightarrow{P}$ là lực hấp dẫn do Trái Đất tác dụng lên một vật được tính bởi công thức $\overrightarrow{P}=mg$, trong đó $m$ là khối lượng của vật (đơn vị: kg), $\overrightarrow{g}$ là vectơ gia tốc rơi tự do, có hướng đi xuống và có độ lớn $g=9,8m/s^{2}$. Xác định độ lớn của trọng lực (đơn vị: N) tác dụng lên quả bóng có khối lượng 450 gam (làm tròn đến hàng phần trăm).

\answer{4,41}

\end{question}

% ===== III-3 (short_answer, difficulty=4) =====
\begin{question}[type=short_answer,difficulty=4,code={III-3}]
Để chặn đường hành lang hình chữ L, người ta dùng một que sào thắng dài đặt kín những điểm chặn với hành lang (như hình vẽ). Biết $ a = 24, b = 3 $. Hỏi cái sào thoả mã điều kiện trên có chiều dài tối thiểu là bao nhiêu? ( kết quả làm tròn đến hàng phần chục).

\image[alt={Hình minh hoạ}]{https://pending-r2.local/raw/21_e450647bab.jpg}


\answer{33,5}

\end{question}

% ===== III-4 (short_answer, difficulty=4) =====
\begin{question}[type=short_answer,difficulty=4,code={III-4}]
Trong cuộc gặp mặt dặn dò trước khi lên đường tham gia kì thi học sinh giỏi, có 10 bạn trong đội tuyển gồm 2 bạn đến từ lớp 12A, 3 bạn từ lớp 12B, 5 bạn còn lại đến từ 5 lớp khác (mỗi lớp một bạn). Thầy giáo xếp ngẫu nhiên các bạn kể trên ngồi vào một bàn dài có 10 ghế mà mỗi bên có 5 ghế xếp đối diện nhau. Xác suất để không có học sinh nào cùng lớp ngồi đối diện nhau bằng $ \frac{a}{b} $ (với $ \frac{a}{b} $ là phân số tối giản). Tính $ a-b $.

\answer{-25}

\end{question}

% ===== III-5 (short_answer, difficulty=4) =====
\begin{question}[type=short_answer,difficulty=4,code={III-5}]
Mảnh đất vườn của nhà anh Điệp có một phần ranh giới cũng là một phần đường cong (C):
 $ y=\frac{x+a}{x+b} $, bao quanh nó là sông nước. Với hệ trục tọa độ Oxy, thích hợp, đơn vị trên mỗi trục là
10 mét thì đường cong (C) đi qua điểm (2; 3) và có đường tiệm cận đứng x=1. Hàng ngày anh Điệp phải dùng thuyền máy để vận chuyển trái cây từ khu vườn của mình đến hai tuyến đường
 $ \Delta_{1}:2x+y-4=0 $ và $ \Delta_{2}:x+2y-2=0 $ cho những người lái buôn từ nơi khác đến. Anh Điệp cần xác định một vị trí $ M(x_{0};y_{0}) $ thuộc khu vườn của mình để tổng các khoảng cách từ vị trí M đó đến hai tuyến đường $ \Delta_{1} $, $ \Delta_{2} $ là bé nhất. Hỏi khoảng cách từ vị trí được chọn làm gốc tọa độ đến điểm M là bao nhiêu mét (làm tròn đến hàng phần chục)?

\image[alt={Hình minh hoạ}]{https://pending-r2.local/raw/25_d1a223ac8a.jpg}


\answer{34,1}

\end{question}

% ===== III-6 (short_answer, difficulty=1) =====
\begin{question}[type=short_answer,difficulty=1,code={III-6}]
Cho hàm số $y = \frac{2x + 3}{1 - x}$. Gọi đường thẳng tiệm cận đứng và tiệm cận ngang của đồ thị hàm số lần lượt có phương trình: $x = a$; $y = b$. Khi đó tổng $a + 3b$ bằng bao nhiêu?

\answer{-5}

\end{question}